\typeout{NT FILE chapter6.tex}%

\chapter{OCamlFLAT Library Implementation}
\label{cha:transducer-implementation}

This chapter details the implementation of finite-state transducer (FST) support within the OCamlFLAT library. 
All logic described here resides primarily in two source files: 
\texttt{TransducerSupport.ml}, which defines the core data types, 
and \texttt{Transducer.ml}, which implements all algorithmic functionality. 
The implementation follows the same architectural conventions as the rest of OCamlFLAT, 
namely a clean separation between type definitions, validation, execution, and higher-level algorithms, 
each encapsulated in its own OCaml module.

\section{Transducer Representation}
\label{sec:back_representation}

The foundation of the FST representation in the library is the OCaml record type \texttt{t}, 
defined inside the \texttt{TransducerBasics} submodule of \texttt{TransducerSupport.ml}. 
This record captures every component of the formal definition of a finite-state transducer as a 6-tuple:

\begin{verbatim}
type t = {
  inAlphabet  : symbols;
  outAlphabet : symbols;
  states      : states;
  initialState: state;
  transitions : transitions4;
  acceptStates: states
}
\end{verbatim}

The field \texttt{inAlphabet} holds the set of input symbols and \texttt{outAlphabet} 
holds the set of output symbols. These are kept as separate sets rather than unified, 
which makes it straightforward to validate that transition symbols are drawn from the correct alphabet and to enforce the 
semantic distinction between reading and writing. The \texttt{states}, \texttt{initialState}, and \texttt{acceptStates} 
fields mirror the corresponding components found in finite automata representations.

The most structurally significant type is \texttt{transitions4}, which is defined as a set of 4-tuples:

\begin{verbatim}
type transition4  = state * symbol * symbol * state
type transitions4 = transition4 set
\end{verbatim}

Each element \texttt{(src, inputSym, outputSym, dst)} represents a single transition: reading symbol 
\texttt{inputSym} in state \texttt{src} causes the machine to emit \texttt{outputSym} and move to state 
\texttt{dst}. Both \texttt{inputSym} and \texttt{outputSym} may take the value \texttt{epsilon} 
(represented as the string \texttt{\textasciitilde} in JSON and as the OCaml constant \texttt{epsilon} internally), 
permitting $\varepsilon$-transitions on the input side and silent transitions on the output side. 
This 4-tuple representation is the key distinction from the 3-tuple \texttt{(src, sym, dst)} used for finite automata, 
and it is consistently used across all modules that manipulate transducers.

The execution engine additionally defines a dedicated configuration type to track the full state of a computation at each step:

\begin{verbatim}
type configuration = state * word * word
type configurations = configuration set
type path           = configuration list
type trail          = configurations list
\end{verbatim}

A \texttt{configuration} is a triple \texttt{(currentState, remainingInput, accumulatedOutput)}, 
representing the current control state, the portion of the input word not yet consumed, and the sequence of output symbols 
produced so far. A \texttt{path} is an ordered list of configurations tracing a single accepting run, and a \texttt{trail} 
collects the full set of configuration sets visited at each step of the computation, which is used to drive the step-by-step 
visualisation in the frontend.

Serialisation and deserialisation are handled by the \texttt{TransducerConversions} module. The \texttt{fromJSon} 
function reads a JSON object and constructs the OCaml record, using helper functions such as \texttt{JSon.fieldSymbolSet}, 
\texttt{JSon.fieldStateSet}, and \texttt{JSon.fieldQuadsSet} to parse each field. The corresponding \texttt{toJSon} and 
\texttt{toJSon2} functions produce JSON representations suitable for saving, loading, and transmitting the transducer definition. 
The \texttt{kind} field is set to the string \texttt{"transducer"}, which is the key used throughout the system to dispatch 
JSON objects to the correct model parser.

\section{Validation and Properties}
\label{sec:back_validation}

Before any algorithm is applied to a transducer, the \texttt{validate} function in \texttt{TransducerPrivate} performs a series 
of structural integrity checks. It verifies that neither alphabet contains epsilon, that the initial state is a member of the 
state set, that all accept states belong to the state set, and that all states and symbols appearing in transitions are drawn 
from their respective declared sets. Each failing check raises a descriptive error through the library's \texttt{Error} module. 
This validation is invoked automatically whenever a transducer is constructed via the \texttt{make} entry point.

Beyond basic structural validation, the implementation provides a rich set of property-checking functions that are exposed to the 
exercise system and to the frontend.

\paragraph{Determinism.}
The function \texttt{isDeterministic} checks whether the transducer is deterministic in the strict sense,
meaning it has no $\varepsilon$-input transitions and each (state, input symbol) pair leads to at most one transition 
with at most one output symbol.

The function proceeds in two stages. First, it performs a global scan over all transitions to test for the presence 
of any $\varepsilon$-input transition, with the form \texttt{(\_,~epsilon,~\_,~\_)}:

\begin{verbatim}
let has_epsilon =
  Set.exists (fun (_, b, _, _) -> b = epsilon) fst.transitions
in
if has_epsilon then false
\end{verbatim}

If any such transition exists the function immediately returns \texttt{false}, since $\varepsilon$-input transitions 
are inherently nondeterministic in this strict model: the machine can advance its state without consuming any input, 
which means that for the same input position there may be multiple reachable states.

If no $\varepsilon$-transitions are present, the function proceeds to the second stage: a nested iteration over every 
state and every input symbol in the declared input alphabet. For each pair \texttt{(st, input)}, it filters the 
full transition set to collect all transitions whose source state is \texttt{st} and whose input symbol is 
\texttt{input}. It then extracts two derived sets, the set of output symbols and the set of destination states 
and requires that all three quantities (number of matching transitions, number of distinct outputs, number of 
distinct destinations) are at most one:

\begin{verbatim}
Set.for_all (fun st ->
  Set.for_all (fun input ->
    let trs  = Set.filter (fun (a,b,_,_) -> a = st && b = input) fst.transitions in
    let outs  = Set.map (fun (_,_,out,_) -> out)  trs in
    let dests = Set.map (fun (_,_,_,d)   -> d)    trs in
    Set.size trs <= 1 && Set.size outs <= 1 && Set.size dests <= 1
  ) fst.inAlphabet
) fst.states
\end{verbatim}

The three conditions together enforce the full determinism requirement. \texttt{Set.size trs <= 1} ensures there is 
at most one applicable transition; \texttt{Set.size outs <= 1} ensures that transition has a unique output symbol; 
and \texttt{Set.size dests <= 1} ensures it leads to a unique next state. All three must hold simultaneously because 
the sets are derived independently and it is theoretically possible, due to the way transitions are stored, for 
the output set or destination set to have size greater than one even when the transition set has size one, if two 
transitions agree on source and input but differ on output or destination.

The overall result is \texttt{true} only if both stages pass for every state and every input symbol, 
confirming the transducer is fully deterministic with no ambiguity in either state or output.

A more permissive function, \texttt{isDeterministicEpsilon}, was attempted, which handles transducers that contain 
$\varepsilon$-transitions but may still be functionally deterministic.

\paragraph{Completeness.}
The function \texttt{isComplete} checks whether for every state and every non-epsilon input symbol there exists at least one 
outgoing transition. This is a standard completeness condition. It is used as a prerequisite for both the Moore and Mealy checks, 
since incomplete machines cannot be classical Mealy or Moore machines by definition.

\paragraph{Moore and Mealy Machines.}
The functions \texttt{isMooreMachine} and \texttt{isMealyMachine} verify membership in the two classical subclasses of FSTs.

For Moore machines, the key invariant is that the output symbol associated with a transition depends only on the source state, 
not on the input symbol. The implementation checks this by examining, for each state, the set of output symbols appearing on 
all outgoing transitions. If this set has cardinality at most one, the state's output is uniform and the Moore condition holds:

\begin{verbatim}
let state_has_single_output =
  Set.for_all (fun st ->
    let outgoing = Set.filter (fun (a,_,_,_) -> a = st) fst.transitions in
    let outputs  = Set.map   (fun (_,_,c,_)  -> c)     outgoing in
    Set.size outputs <= 1
  ) fst.states
in
deterministic && complete && state_has_single_output
\end{verbatim}

For Mealy machines, the condition is weaker: the output must be uniquely determined by the pair (state, input symbol), 
rather than by state alone. The implementation checks that for every (state, input) pair, the set of output symbols across 
all matching transitions has cardinality at most one. Both functions also require the machine to be deterministic and complete.

\paragraph{Reachability, Productivity, and Cleanliness.}
The library also implements state-level analyses mirroring those of finite automata. The function \texttt{reachable} 
computes the set of all states reachable from a given state via any sequence of transitions, using iterative forward exploration. 
The function \texttt{productive} identifies states from which some accepting state can be reached. 
The function \texttt{getUsefulStates} returns the intersection of reachable and productive states, 
and \texttt{cleanUselessStates} constructs an equivalent transducer containing only useful states and 
only the transitions between them, recomputing the alphabets from the surviving transitions. The function \texttt{isClean} 
delegates to the corresponding finite automaton analysis by projecting the transducer to its underlying automaton structure 
using \texttt{asFiniteAutomaton}, which strips the output symbols from transitions.

\section{Execution Engine}
\label{sec:back_execution}

The execution engine is implemented in the \texttt{TransducerAccept} module and follows the same general framework used for 
finite automata and pushdown automata throughout OCamlFLAT. Computation proceeds by maintaining a set of active configurations 
and expanding them step by step through the \texttt{nextConfigs} function.

The initial configuration is constructed by \texttt{initialConfigs}, which wraps the start state, the full input word, 
and an empty output accumulator into a singleton set:

\begin{verbatim}
let initialConfigs (fst: t) (w: word) : configurations =
  Set.make [(fst.initialState, w, [])]
\end{verbatim}

The \texttt{nextConfigs} function defines the single-step transition relation over configurations. 
Given a configuration \texttt{(st, w, out)}, it considers two cases. If the remaining input \texttt{w} is empty, 
only $\varepsilon$-input transitions from \texttt{st} are applicable, and each such transition 
\texttt{(st, epsilon, outSym, st2)} produces a new configuration \texttt{(st2, [], out @ [outSym])} 
(with epsilon output symbols suppressed). If \texttt{w = x :: xs}, both symbol-consuming transitions on \texttt{x} 
and $\varepsilon$-input transitions on the current state are applicable, and the results are merged:

\begin{verbatim}
| x::xs ->
  let consume = Set.filter (fun (st1,sy,_,_) -> st1 = st && sy = x) fst.transitions in
  let viaConsume = Set.map (build_next_config xs) consume in
  let epsTr = Set.filter (fun (st1,sy,_,_) -> st1 = st && sy = epsilon) 
  fst.transitions in
  let viaEps = Set.map (build_next_config w) epsTr in
  Set.union viaConsume viaEps
\end{verbatim}

The helper \texttt{build\_next\_config} constructs the next configuration given the new remaining input and the matched transition, 
appending the output symbol to the accumulator when it is not epsilon.

A configuration is accepting when the control state is in \texttt{acceptStates} and the remaining input word is empty. 
The plain \texttt{accept} function simply tests whether any configuration in the reachable set eventually becomes accepting, 
delegating to the generic \texttt{Model.accept} combinator.

The function \texttt{acceptFull} returns a triple \texttt{(accepted, path, trail)}, where \texttt{path} 
is the sequence of configurations along one accepting run and \texttt{trail} is the full history of all 
configuration sets visited at each step. The trail is what the frontend uses to drive animated, step-by-step execution: 
each element of the trail corresponds to one round of expansion, and the frontend renders the current state, 
remaining input, and accumulated output at each step.

The function \texttt{acceptOut} is a convenience wrapper over \texttt{acceptFull} that extracts only the output word from 
the final configuration of an accepting run, returning a pair \texttt{(true, outputWord)} on acceptance or \texttt{(false, [])} 
on rejection. This is the function used to answer the question ``what does this transducer produce on this input?'', 
as opposed to the binary question of acceptance.

Two additional helpers enforce output validation for exercises. \texttt{acceptCheckOutput} first checks that the input belongs to
the declared input alphabet, runs \texttt{acceptOut}, and then verifies that the produced output word belongs to the declared
output alphabet. \texttt{acceptExpectedOutput} extends this by supporting exercise words written as \texttt{input->output}:
it parses the arrow, feeds the input part to the transducer, and accepts only if the actual output exactly matches the expected
one (and still belongs to the output alphabet). The exercise harness now uses \texttt{acceptExpectedOutput}, so authors can list
cases such as \texttt{a->a}, \texttt{ab->ab}, etc., in the \emph{inside} set to assert both acceptance and the precise output; any
word without an arrow falls back to the alphabet-checked acceptance behaviour.

Output-aware acceptance helper:
\begin{verbatim}
let acceptExpectedOutput (fst: t) (w: word) : bool =
  let wstr = word2str w in
  match String.index_opt wstr '-' with
  | Some i when i + 1 < String.length wstr && wstr.[i + 1] = '>' ->
      let inp = String.sub wstr 0 i |> str2word in
      let out_exp = String.sub wstr (i + 2) (String.length wstr - i - 2) |> str2word in
      let (ok, out_act) = acceptOut fst inp in
      ok && out_act = out_exp && Model.checkWord fst.outAlphabet out_act
  | _ -> acceptCheckOutput fst w
\end{verbatim}
This helper matches the library implementation: parse the arrow form, run the transducer on the input, 
and only accept when the produced output matches exactly and belongs to the declared output alphabet.

The \texttt{TransducerGenerate} module provides word generation functionality. Its \texttt{nextConfigs2} 
function reverses the direction of exploration: rather than consuming an input word, it explores outgoing 
transitions from the current state freely, accumulating both the input symbols read and the output symbols emitted. 
The \texttt{generate} function uses this to enumerate all input words of a given length that the transducer accepts, 
relying on the generic \texttt{Model.generate} infrastructure.

\section{Algorithms}
\label{sec:back_algorithms}

\subsection{Determinization}
The determinization algorithm is one of the most complex components of the FST implementation. 
%Two variants were implemented, but ultimately only the first one is used.

The implementation choices made for this algorithm were made after research and implementation attempts, and were necessitated by fundamental system limitations 
within the OCamlFLAT/OFLAT architecture. Specifically, the system's transition data structure enforces that an output symbol 
on a transition must be a single symbol rather than an arbitrary output string or sequence. 
Consequently, supporting general delayed-output determinization, where $\varepsilon$-transitions 
carry outputs and varying output lengths accumulate across subset paths would violate this single-symbol transition constraint 
and break UI rendering and table mappings. While it could have been done, it would require significant changes to parts of the library that
didn't pertain to me. So the decision to restrict determinization to not handling transducers with output-producing $\varepsilon$-transitions,
was made, as it bypasses the mentioned constraints, ensuring predictable execution, and a reliable interface for the web application.

\texttt{toDeterministic}, handles all transducers except those with output-producing $\varepsilon$-transitions. 
It begins by rejecting the input if any $\varepsilon$-transition has a non-epsilon output symbol. 
If this check passes, it applies a standard subset construction used in finite automata, adapted for transducers. 
The algorithm maintains a worklist of sets of NFST states (each representing a single DFST state), 
starts from the $\varepsilon$-closure of the initial state, and iterates by computing for each unprocessed state 
set and each input symbol the set of states reachable via that symbol followed by $\varepsilon$-closure. 
If at any point two transitions with the same source state and input symbol are found to produce different output symbols, 
the algorithm raises an internal exception \texttt{DeterminizationFailed} and returns the original transducer unchanged. 
The state names in the resulting deterministic transducer are synthesised by the \texttt{fusedfstState} function, 
which concatenates the names of the constituent NFST states into a canonical string representation enclosed in curly braces.



%The other variant, \texttt{toDeterministicEpsilon}, handles transducers with output-producing $\varepsilon$-transitions. 
%This requires a more subtle algorithm based on the concept of \emph{delayed output}. Each state in the determinized transducer 
%is represented not as a set of NFA states, but as a set of \texttt{(state, word)} pairs, where the word component tracks the 
%output that has been accumulated along $\varepsilon$-transitions but not yet committed to a transition label. The algorithm is:

%\begin{enumerate}
%  \item Begin with the $\varepsilon$-output-closure of the initial state, computed by \texttt{closeEmptyOut}.
%  \item For each current DFST state (a set of \texttt{(state, word)} pairs) and each input symbol, compute all reachable 
%  \texttt{(nextState, fullOutputWord)} pairs by following symbol-consuming transitions and then taking $\varepsilon$-closures.
%  \item Compute the longest common prefix (LCP) of all output words in the result. This LCP is the output emitted on the transition. 
%  Subtract the LCP from each word to form the delayed suffix, and use the resulting \texttt{(state, suffix)} pairs as the next DFST 
%  state.
%  \item If the LCP is non-empty, its first symbol becomes the output label of the new transition. Otherwise, the transition is silent.
%  \item Continue until no new DFST states are discovered.
%\end{enumerate}

%The LCP computation is handled by the \texttt{lcp} and \texttt{suffix} helper functions, and the consistency of outputs at each step
% is enforced by the \texttt{prefix\_consistent} and \texttt{hasConflictOut} checks. If the algorithm cannot commit to a 
% single output prefix at any step, it reports an error and returns the original transducer.

%Delayed-output determinization step:
%\begin{verbatim}
%let getNext (q: dfstState) (sy: symbol) (all_transitions: transitions4) =
%  let allResults =
%    Set.flatMap
%      (fun (st, outPrefix) ->
%        let moveTransitions =
%          Set.filter (fun (a, b, _, _) -> a = st && b = sy) all_transitions
%        in
%        Set.flatMap
%          (fun (_, _, moveOut, destSt) ->
%            let newPrefix = outPrefix @ [moveOut] in
%            let epsClosure = closeEmptyOut destSt all_transitions in
%            Set.map (fun (s, epsOut) -> (s, newPrefix @ epsOut)) epsClosure)
%          moveTransitions)
%      q
%  in
%  if Set.isEmpty allResults then None
%  else
%    let allWords = Set.map (fun (_, w) -> w) allResults in
%    let firstWord = (Set.toList allWords) |> List.hd in
%    let lcpWord = Set.fold_left lcp firstWord allWords in
%    let newdfstState = Set.map (fun (s, w) -> (s, suffix lcpWord w)) allResults in
%    let outSymbol = if List.length lcpWord > 0 then List.hd lcpWord else epsilon in
%    Some (outSymbol, newdfstState)
%\end{verbatim}
%The determinization step computes reachable \texttt{(state, word)} pairs, commits the longest common prefix as the emitted
%output, and propagates any remaining suffixes as delayed output for the successor DFST state.


\subsection{Minimization}

The \texttt{minimize} function reduces a deterministic transducer to its minimal equivalent using a partition-refinement 
algorithm analogous to Hopcroft's algorithm for finite automata~\cite{hopcroft2006automata}, but adapted to account for output symbols,
following the approach described by Mohri~\cite{mohri1998minimization}.

The algorithm begins by removing all useless states via \texttt{cleanUselessStates}, then partitions the remaining states into 
two initial blocks: accepting states and non-accepting states. It then iteratively refines this partition. 
In each refinement step, every block is split according to the \emph{signature} of its member states, 
where the signature of a state is the set of triples \texttt{(inputSymbol, outputSymbol, destinationBlock)} 
derived from its outgoing transitions. Two states have the same signature if and only if they are behaviourally 
indistinguishable with respect to input consumption, output production, and the equivalence class of their successors. 
The refinement continues until no further splitting is possible:

\begin{verbatim}
let rec refine (partition: states Set.t) : states Set.t =
  let refined = Set.flatMap
    (fun block -> Set.make (refine_block partition block)) partition in
  if Set.equals refined partition then partition
  else refine refined
\end{verbatim}

Once the fixed point is reached, a representative state is chosen from each equivalence class 
(the first element of the class, as ordered by the set representation), and all transitions are re-mapped to use representatives. 
The resulting transducer has one state per equivalence class and is guaranteed to be the unique minimal deterministic 
transducer recognising the same input-output relation.

The function \texttt{isMinimized} checks whether a deterministic transducer is already minimal by running 
\texttt{minimize} and comparing the number of states. If the count is unchanged, the transducer is already minimal.

\subsection{Conversion to Turing Machine}

The function \texttt{asTuringMachine} converts an FST into an equivalent 2-tape Turing machine, 
where tape 1 acts as a read-only input tape and tape 2 acts as a write-only output tape. 
This is used to demonstrate the relationship between transducers and more general models of computation.

The conversion works as follows. Each FST transition \texttt{(q1, a, b, q2)} where \texttt{a} is not epsilon becomes
 a TM transition that reads \texttt{a} from tape 1, reads the blank symbol from tape 2, 
 moves to state \texttt{q2}, writes \texttt{a} back on tape 1 (keeping it read-only), writes \texttt{b} on tape 2, 
 and advances both heads rightward. If the output symbol \texttt{b} is epsilon, the tape 2 head stays in place and 
 the blank is preserved. For $\varepsilon$-input transitions, a separate set of TM transitions is generated for 
 each possible symbol on tape 1, keeping the tape 1 head stationary while still writing to tape 2. Finally, 
 a dedicated accepting state \texttt{q\_accept\_tm} is added, and each FST accepting state receives a transition 
 to it when both tapes read the blank symbol, signalling that the input has been fully consumed and processing is complete.

\section{Composition and Operations}
\label{sec:back_operations}

The \texttt{TransducerComposition} module implements five structural operations on transducers: composition, 
union, intersection, inverse, and concatenation. All are implemented using worklist-based state-space exploration 
or simple structural transformations.

\paragraph{Composition.}
Given two transducers $T_1 : A \to B$ and $T_2 : B \to C$, their composition $T_1 \circ T_2 : A \to C$ is the 
transducer that feeds the output of $T_1$ directly into $T_2$. The implementation uses a product-construction 
worklist over pairs of states \texttt{(q1, q2)}, one from each transducer. Three cases of synchronisation are handled:

\begin{itemize}
  \item \textbf{Lock-step:} $T_1$ consumes a non-epsilon input symbol \texttt{a} and emits 
  \texttt{b}; $T_2$ simultaneously consumes \texttt{b} and emits \texttt{c}. 
  The composed transition is \texttt{(q1,q2) --a/c--> (q1',q2')}.
  \item \textbf{$T_1$ epsilon-step:} $T_1$ takes an $\varepsilon$-input transition emitting 
  \texttt{b}; $T_2$ consumes \texttt{b} and emits \texttt{c}. The composed transition is labelled with $\varepsilon$ input.
  \item \textbf{$T_2$ epsilon-step:} $T_2$ takes an $\varepsilon$-input transition independently; 
  $T_1$ stays at its current state. The composed transition is a pure $\varepsilon$ step.
\end{itemize}

State names in the composed transducer are formed by concatenating the names of the component states with an 
underscore separator via the \texttt{fuse\_states} helper. A pair of states is accepting in the composition 
if and only if both component states are accepting.

\paragraph{Union.}
The union $T_1 \cup T_2$ is constructed by introducing a fresh initial state and adding two silent 
$\varepsilon$-transitions, one to the renamed initial state of $T_1$ and one to the renamed initial state of $T_2$. 
The states and transitions of both machines are renamed with suffix \texttt{\_1} and \texttt{\_2} respectively to 
avoid name collisions, and their sets are merged. This is the standard nondeterministic union construction for automata, 
directly adapted to transducers.

\paragraph{Intersection.}
The intersection $T_1 \cap T_2$ is defined as the transducer whose accepted pairs \texttt{(w, v)} are exactly 
those accepted by both $T_1$ and $T_2$. The implementation uses a product construction where a transition 
is added to the result only when both machines simultaneously make a transition with identical input and identical output. 
This is more restrictive than composition and captures the set-theoretic intersection of the input-output relations. 
Three synchronisation cases are handled: lock-step moves where both machines agree on input and output, and two cases 
for independent silent $(\varepsilon, \varepsilon)$ steps by either machine.

\paragraph{Inverse.}
The inverse $T^{-1}$ of a transducer $T$ is formed by swapping the roles of the input and output alphabets and 
swapping the input and output symbols in every transition. The implementation is a simple structural map over the 
transition set and a swap of \texttt{inAlphabet} and \texttt{outAlphabet}. If $T$ maps input word $w$ to output word $v$, 
then $T^{-1}$ maps $v$ to $w$.

\paragraph{Concatenation.}
The concatenation $T_1 \cdot T_2$ accepts a pair \texttt{(w, v)} if and only if there exist $w = w_1 w_2$ and 
$v = v_1 v_2$ such that $T_1$ accepts \texttt{(w1, v1)} and $T_2$ accepts \texttt{(w2, v2)}. 
The implementation renames both machines to avoid state collisions, connects every accepting state of $T_1$ 
to the initial state of $T_2$ via a silent $(\varepsilon, \varepsilon)$ transition, and sets the accepting states of the result 
to be those of $T_2$ alone. The initial state of the concatenation is that of $T_1$.

\section{Integration Points Exposed to OFLAT}

The library surface made available to the OFLAT frontend aligns with the features highlighted earlier:
\begin{itemize}
  \item \textbf{Acceptance with outputs}: \texttt{acceptOut} returns the produced output word; 
  \texttt{acceptExpectedOutput} enforces equality against an expected word (parsed from \texttt{input->output}) and validates both alphabets. 
  These power output-aware exercises and quick checks.
  \item \textbf{Stepwise visualization}: \texttt{acceptFull} returns both a sample accepting path and the full trail of configuration sets. 
  The frontend consumes the trail to highlight active states/transitions (via \texttt{inputEdges}) 
  and to display the evolving output during stepping.
  \item \textbf{Generation}: \texttt{generate} enumerates accepted input/output pairs of bounded length, reused by the generic 
  generation UI and exercise scaffolding.
  \item \textbf{Transformations}: determinization/minimization, inverse, union/intersection/concatenation, 
  and composition are exposed so OFLAT can chain models and present derived transducers directly to the user.
  \item \textbf{JSON interoperability}: \texttt{toJSon}/\texttt{fromJSon} keep a stable \texttt{kind = "transducer"} tag, 
  ensuring saved files load correctly across the web client and the OCaml library.
\end{itemize}

These hooks ensure the theoretical and algorithmic support described above is consumable by the web UI without duplicating logic in JavaScript.

